%=============================================================================!
% 8.8  LONG-RANGE ELECTROSTATICS                                              !
%=============================================================================!
\section{Charges and the slow decay of \texorpdfstring{$1/r$}{1/r}}
\label{sec:pe-coulomb}

Ions in an electrolyte, the charged atoms of a crystal such as rock salt
and the charges placed on the atoms of a force field interact by
Coulomb's law. Two charges
$q_i$ and $q_j$, in units of the elementary charge $e$, the charge of a
proton (an electron carries $-e$, and an ion is an atom that has lost or
gained electrons), at a distance $r$ have the energy
\begin{equation}
   \varphi_{\mathrm C}(r) = \frac{k_{\mathrm e}\,q_iq_j}{r} ,
   \qquad k_{\mathrm e} = \frac{e^2}{4\pi\epsilon_0} = \SI{14.3996}{\electronvolt
   \angstrom} ,
   \label{eq:pe-coulomb}
\end{equation}
where $\epsilon_0$ is the electric constant; $k_{\mathrm e}$ is computed from its
value in \texttt{mdlab.units}. This section shows that this energy cannot be
cut off as the earlier ones were.

\subsection*{A slow decay}

Two opposite unit charges at \SI{1.5639}{\angstrom}, the bond length of
LiF (\cref{sec:ro-molecules}), have $-14.3996/1.5639 =
\SI{-9.21}{\electronvolt}$; at \SI{10}{\angstrom} they still have
$\SI{-1.44}{\electronvolt}$, where two argon atoms have
$\SI{-6.4e-5}{\electronvolt}$. Take an ion inside a large region filled with ions at the number density
$n$, that is $n$ ions in each unit of volume. The ions at distances between $r$ and $r + \Delta r$ fill a
shell of area $4\pi r^2$ and thickness $\Delta r$, so there are about
$4\pi r^2n\,\Delta r$ of them, a number that grows as $r^2$. Each
contributes an energy of size $k_{\mathrm e}/r$, so the shell as a whole contributes
an amount of size $4\pi k_{\mathrm e}n\,r\,\Delta r$, which grows with $r$. For a
Lennard-Jones attraction the same shell contributes $4\pi r^2n\,\Delta
r\times r^{-6}$, falling as $r^{-4}$, and the far shells hardly matter;
for charges, every shell counts more than the last, and only the
cancellation between positive and negative ions keeps the total finite.

\subsection*{A sum that depends on its order}

Rock salt, NaCl, shows what this means. Its ions sit on a cubic grid of
spacing $d$, the sign of the charge alternating from each site to its
neighbours, so the ion at $(a, b, c)d$, with $a$, $b$, $c$ integers, positive,
negative or zero,
carries the charge $(-1)^{a+b+c}$ relative to the ion at the origin. The
energy of the ion at the origin with all the others is $-\mathcal{M}k_{\mathrm e}/d$, with the Madelung
constant
\begin{equation*}
   \mathcal{M} = -\sum_{(a,b,c)\ne(0,0,0)}\frac{(-1)^{a+b+c}}{\sqrt{a^2 + b^2 + c^2}} .
\end{equation*}
\Cref{fig:pe-madelung} adds up this sum in two orders. Over the ions
inside a growing sphere, the partial sums do not settle, out to $12d$: the six nearest
ions alone give the partial sum 6, adding the twelve next give $-2.49$, and between
radii $6d$ and $12d$ the sums swing between $-10.3$ and $+14.7$, because
each new shell of ions has a net charge of its own. Over the ions inside a
growing cube, with the ions on its faces, edges and corners counted with
the fractions $\frac12$, $\frac14$ and $\frac18$ of a small cube of side
$d$ around each that lie inside, every cube is neutral: the weighted
charge is a product of three equal sums, one along each axis, each zero,
as $-\frac12 + 1 - \frac12 = 0$ for the smallest cube. The sums settle
quickly, to 1.747567 at a half-side of $12d$, within $2\times10^{-6}$ of
the limit 1.747565\cite{Evjen1932}. The answer depends on the order in
which the terms are added. That cannot happen when the sizes of the
terms, all taken positive, add up to a finite total; here they do not,
since the sizes in one shell add up to an amount that grows as $r$.
Physically, the arrangement of charge at the surface of the
region summed over affects the energy of every ion inside it. For a crystal of this kind with a spacing $d =
\SI{2.82}{\angstrom}$, the electrostatic energy of one ion with all the
others is $-1.7476\times14.3996/2.82 = \SI{-8.92}{\electronvolt}$.

\begin{figure}[htb]
\centering
\includegraphics{ch08_potentials/madelung}
\caption{The Madelung sum of rock salt over the ions inside a growing
sphere (ochre), which never settles, and inside a growing neutral cube
(teal dots), which settles to 1.7476 (dotted), against the radius or
half-side in units of the nearest-neighbour distance $d$.}
\label{fig:pe-madelung}
\end{figure}

A cutoff therefore cannot be applied to charges as it was to the
Lennard-Jones energy: whatever its radius, the charge inside it is not
zero in general and the energy depends on where it falls. Electrolytes,
full of ions, and electrodes, whose surfaces carry charge, need the
electrostatic energy of an infinite, periodic arrangement of charges
summed correctly. Chapter~10, once periodic boxes are in place, derives
the Ewald method, one way of doing that.

\begin{keyidea}
The Coulomb energy $k_{\mathrm e}q_iq_j/r$ falls so slowly that each shell of
charges matters more than the last, and the sum for a crystal of ions
depends on the order in which its terms are added, which sets the charge
at the surface of the region summed. Charges cannot be cut off like a
Lennard-Jones energy, and periodic systems need a sum built for them,
such as the Ewald sum of Chapter~10.
\end{keyidea}

\notebook{08}{sum the charges of rock salt over spheres and cubes}

\begin{selfcheck}
How many ions lie at the distance $\sqrt2\,d$ from an ion in rock salt,
and what is their charge relative to it?
\end{selfcheck}
